\documentclass[../libro]{subfiles}

\begin{document}

\section{Pfaffian 的完全展开}

设 \(A\) 是 \(n\)~级阵.
我们定义了 \(A\) 的行列式
\begin{align*}
    \det {(A)}
    =
    \begin{dcases}
        [A]_{1,1},
         & n = 1;    \\
        \sum_{i = 1}^{n}
        {(-1)^{i+1} [A]_{i,1} \det {(A(i|1))}},
         & n \geq 2.
    \end{dcases}
\end{align*}
然后, 我们证明了, 具体地,
\begin{equation}
    \det {(A)}
    =
    \sum_{\substack{
            1 \leq i_1, i_2, \dots, i_n \leq n; \\
            i_1, i_2, \dots, i_n\,\text{互不相同}
        }}
    {s(i_1, i_2, \dots, i_n)\,
        [A]_{i_1,1} [A]_{i_2,2} \dots [A]_{i_n,n}},
    \label{eq:explicit-expression-of-determinant}
\end{equation}
其中,
\begin{align*}
    s(i_1, i_2, \dots, i_n)
    = {} & \prod_{1 \leq u < v \leq n}
    {\operatorname{sgn} {(i_v - i_u)}},
\end{align*}
其中,
\begin{align*}
    \operatorname{sgn} {(x)}
    =
    \begin{cases}
        1,  & x > 0; \\
        0,  & x = 0; \\
        -1, & x < 0.
    \end{cases}
\end{align*}
式~\eqref{eq:explicit-expression-of-determinant}
是行列式的完全展开.

设 \(A\) 是 \(n\)~级反称阵.
我们定义了 \(A\) 的 pfaffian
\begin{align*}
    \operatorname{pf} {(A)} =
    \begin{dcases}
        0,
         & n = 1;    \\
        [A]_{1,2},
         & n = 2;    \\
        \sum_{j = 2}^{n}
        {(-1)^{j} [A]_{1,j}
        \operatorname{pf} {(A({1,j}|{1,j}))}},
         & n \geq 3.
    \end{dcases}
\end{align*}
当 \(n\) 是奇数时,
我们用数学归纳法证明了,
\(\operatorname{pf} {(A)} = 0\).
当 \(n = 2m\) 是偶数时,
我们也能具体地写 \(\operatorname{pf} {(A)}\) 吗?

本节, 我们研究偶数级反称阵的 pfaffian 的完全展开.

\begin{definition}
    设 \(A\) 是 \(2m\)~级反称阵.
    定义 \(A\) 的 ``qfaffian''
    \begin{align*}
        \operatorname{qf} {(A)}
        = \sum_{\substack{
        1 \leq i_1, j_1, \dots, i_m, j_m \leq 2m; \\
        i_1, j_1, \dots, i_m, j_m\,\text{互不相同};   \\
        i_1 < j_1, \dots, i_m < j_m;              \\
                i_1 < \dots < i_m
            }}
        {s(i_1, j_1, \dots, i_m, j_m)\,
            [A]_{i_1,j_1} \dots [A]_{i_m,j_m}}.
    \end{align*}
\end{definition}

注意, 既然 \(i_1 < j_1\), \(\dots\), \(i_m < j_m\),
且 \(i_1 < \dots < i_m\),
则 \(i_1\) 一定比 \(i_2\), \(\dots\), \(i_m\) 都小,
且 \(i_1\) 也比 \(j_1\), \(\dots\), \(j_m\) 都小.
则 \(i_1\) 是 \(i_1\), \(j_1\), \(\dots\), \(i_m\), \(j_m\)
的最小的数.

注意, 我们不定义奇数级反称阵的 ``qfaffian''.

我们计算小级反称阵的 ``qfaffian''.

\begin{example}
    设 \(A\) 是 \(2 \cdot 1\)~级反称阵.
    则, 适合 \(i_1 < j_1\),
    且 \(1 \leq i_1, j_1 \leq 2\) 的 \((i_1, j_1)\) 是 \((1, 2)\).
    则 \(\operatorname{qf} {(A)} = s(1, 2)\, [A]_{1,2} = [A]_{1,2}\).
    注意, \(\operatorname{qf} {(A)} = \operatorname{pf} {(A)}\).
\end{example}

\begin{example}
    设 \(A\) 是 \(2 \cdot 2\)~级反称阵.
    则, 适合 \(i_1 < j_1\), \(i_2 < j_2\),
    \(i_1 < i_2\),
    且 \(1 \leq i_1, j_1, i_2, j_2 \leq 4\)
    的 \((i_1, j_1, i_2, j_2)\) 是
    \((1, 2, 3, 4)\), \((1, 3, 2, 4)\), \((1, 4, 2, 3)\).
    则
    \begin{align*}
             &
        \operatorname{qf} {(A)}
        \\
        = {} &
        s(1, 2, 3, 4)\, [A]_{1,2} [A]_{3,4}
        + s(1, 3, 2, 4)\, [A]_{1,3} [A]_{2,4}
        + s(1, 4, 2, 3)\, [A]_{1,4} [A]_{2,3}
        \\
        = {} &
        1\, [A]_{1,2} [A]_{3,4}
        + (-1)\, [A]_{1,3} [A]_{2,4}
        + 1\, [A]_{1,4} [A]_{2,3}
        \\
        = {} &
        [A]_{1,2} [A]_{3,4}
        - [A]_{1,3} [A]_{2,4}
        + [A]_{1,4} [A]_{2,3}.
    \end{align*}
    注意, \(\operatorname{qf} {(A)} = \operatorname{pf} {(A)}\).
\end{example}

我们会证明, 对任何 \(2m\)~级反称阵 \(A\),
\(\operatorname{qf} {(A)} = \operatorname{pf} {(A)}\).
不过, 我们要作准备.

\begin{theorem}
    设 \(A\) 是 \(2m\)~级反称阵.
    设 \(2 \leq 2k \leq 2m\).
    设 \(1 \leq h_1 < \dots < h_{2k} \leq 2m\).
    则
    \begin{align*}
        \operatorname{qf} {\left(
            A\binom{h_1,\dots,h_{2k}}{h_1,\dots,h_{2k}}
            \right)}
        = \sum_{\substack{
        i_1, j_1, \dots, i_k, j_k \text{是} \\
        h_1, \dots h_{2k} \text{的排列};      \\
        i_1 < j_1, \dots, i_k < j_k;       \\
                i_1 < \dots < i_k
            }}
        {s(i_1, j_1, \dots, i_k, j_k)\,
            [A]_{i_1,j_1} \dots [A]_{i_k,j_k}}.
    \end{align*}
\end{theorem}

\begin{proof}
    记
    \begin{align*}
        B = A\binom{h_1,\dots,h_{2k}}{h_1,\dots,h_{2k}}.
    \end{align*}
    则 \(B\) 是 \(2k\)~级反称阵,
    且 \([B]_{u,v} = [A]_{h_u,h_v}\).

    记 \(f(h_t) = t\), \(t = 1\), \(\dots\), \(2k\).
    则若 \(h_u < h_v\), 必 \(f(h_u) < f(h_v)\);
    反过来, 若 \(f(h_u) < f(h_v)\), 必 \(h_u < h_v\).
    那么:

    (1)
    对任何 \(h_1\), \(\dots\), \(h_{2k}\) 的排列
    \(i_1\), \(j_1\), \(\dots\), \(i_k\), \(j_k\),
    我们有,
    \(f(i_1)\), \(f(j_1)\), \(\dots\), \(f(i_k)\), \(f(j_k)\)
    一定是 \(1\), \(\dots\), \(2k\) 的排列.
    反过来, 对任何 \(1\), \(\dots\), \(2k\) 的排列
    \(u_1\), \(v_1\), \(\dots\), \(u_k\), \(v_k\),
    我们有,
    \(h_{u_1}\), \(h_{v_1}\), \(\dots\), \(h_{u_k}\), \(h_{v_k}\),
    一定是 \(h_1\), \(\dots\), \(h_{2k}\) 的排列.

    (2)
    若 \(h_1\), \(\dots\), \(h_{2k}\) 的排列
    \(i_1\), \(j_1\), \(\dots\), \(i_k\), \(j_k\)
    适合
    \(i_1 < j_1\), \(\dots\), \(i_k < j_k\),
    则对应的 \(1\), \(\dots\), \(2k\) 的排列
    \(f(i_1)\), \(f(j_1)\), \(\dots\), \(f(i_k)\), \(f(j_k)\)
    适合
    \(f(i_1) < f(j_1)\), \(\dots\), \(f(i_k) < f(j_k)\).
    反过来, 若 \(1\), \(\dots\), \(2k\) 的排列
    \(u_1\), \(v_1\), \(\dots\), \(u_k\), \(v_k\)
    适合
    \(u_1 < v_1\), \(\dots\), \(u_k < v_k\),
    则对应的 \(h_1\), \(\dots\), \(h_{2k}\) 的排列
    \(h_{u_1}\), \(h_{v_1}\), \(\dots\), \(h_{u_k}\), \(h_{v_k}\)
    适合
    \(h_{u_1} < h_{v_1}\), \(\dots\), \(h_{u_k} < h_{v_k}\).

    (3)
    若 \(h_1\), \(\dots\), \(h_{2k}\) 的排列
    \(i_1\), \(j_1\), \(\dots\), \(i_k\), \(j_k\)
    适合
    \(i_1 < \dots < i_k\),
    则对应的 \(1\), \(\dots\), \(2k\) 的排列
    \(f(i_1)\), \(f(j_1)\), \(\dots\), \(f(i_k)\), \(f(j_k)\)
    适合
    \(f(i_1) < \dots < f(i_k)\).
    反过来, 若 \(1\), \(\dots\), \(2k\) 的排列
    \(u_1\), \(v_1\), \(\dots\), \(u_k\), \(v_k\)
    适合
    \(u_1 < \dots < u_k\),
    则对应的 \(h_1\), \(\dots\), \(h_{2k}\) 的排列
    \(h_{u_1}\), \(h_{v_1}\), \(\dots\), \(h_{u_k}\), \(h_{v_k}\)
    适合
    \(h_{u_1} < \dots < h_{u_k}\).

    (4)
    既然 \(h_u < h_v\) 相当于 \(f(h_u) < f(h_v)\),
    我们有,
    \(\operatorname{sgn} {(f(h_t) - f(h_s))}
    = \operatorname{sgn} {(h_t - h_s)}\).
    则
    \(s(f(i_1), f(j_1), \dots, f(i_k), f(j_k))
    = s(i_1, j_1, \dots, i_k, j_k)\).

    综上,
    \begin{align*}
             &
        \operatorname{qf} {\left(
            A\binom{h_1,\dots,h_{2k}}{h_1,\dots,h_{2k}}
            \right)}
        \\
        = {} &
        \operatorname{qf} {(B)}
        \\
        = {} &
        \sum_{\substack{
        i_1, j_1, \dots, i_k, j_k \text{是} \\
        h_1, \dots h_{2k} \text{的排列};      \\
        i_1 < j_1, \dots, i_k < j_k;       \\
                i_1 < \dots < i_k
            }}
        {s(f(i_1), f(j_1), \dots, f(i_k), f(j_k))\,
            [B]_{f(i_1),f(j_1)} \dots [B]_{f(i_k),f(j_k)}}
        \\
        = {} &
        \sum_{\substack{
        i_1, j_1, \dots, i_k, j_k \text{是} \\
        h_1, \dots h_{2k} \text{的排列};      \\
        i_1 < j_1, \dots, i_k < j_k;       \\
                i_1 < \dots < i_k
            }}
        {s(i_1, j_1, \dots, i_k, j_k)\,
            [A]_{i_1,j_1} \dots [A]_{i_k,j_k}}.
        \qedhere
    \end{align*}
\end{proof}

\begin{theorem}
    设 \(p_1\), \(p_2\), \(p_3\), \(\dots\), \(p_n\)
    是不超过 \(n\) 的正整数,
    且互不相同.
    设 \(p_1 = 1\).
    则
    \begin{align*}
        s(p_1, p_2, p_3, \dots, p_n)
        = (-1)^{p_2 - 2}\, s(p_3, \dots, p_n).
    \end{align*}
\end{theorem}

\begin{proof}
    注意,
    \begin{align*}
             &
        s(p_1, p_2, p_3, \dots, p_n)
        \\
        = {} &
        \prod_{1 \leq u < v \leq n}
        {\operatorname{sgn} {(p_v - p_u)}}
        \\
        = {} &
        \left(
        \prod_{\substack{1 \leq u < v \leq n \\
                u = 1}}
        {\operatorname{sgn} {(p_v - p_u)}}
        \right)
        \cdot
        \left(
        \prod_{\substack{1 \leq u < v \leq n \\
                u = 2}}
        {\operatorname{sgn} {(p_v - p_u)}}
        \right)
        \cdot
        \left(
        \prod_{\substack{1 \leq u < v \leq n \\
                u \geq 3}}
        {\operatorname{sgn} {(p_v - p_u)}}
        \right)
        \\
        = {} &
        \left(
        \prod_{2 \leq v \leq n}
        {\operatorname{sgn} {(p_v - p_1)}}
        \right)
        \cdot
        \left(
        \prod_{3 \leq v \leq n}
        {\operatorname{sgn} {(p_v - p_2)}}
        \right)
        \cdot
        s(p_3, \dots, p_n).
    \end{align*}
    注意, 既然 \(p_1 = 1\),
    则 \(v \neq 1\) 时, 必 \(p_v > 1\).
    则
    \begin{align*}
        \prod_{2 \leq v \leq n}
        {\operatorname{sgn} {(p_v - p_1)}}
        = 1.
    \end{align*}
    注意, 在 \(p_3\), \(\dots\), \(p_n\) 里,
    恰有 \(p_2 - 2\) 个数比 \(p_2\) 小
    (注意, \(p_1 = 1\)).
    则
    \begin{align*}
        \prod_{3 \leq v \leq n}
        {\operatorname{sgn} {(p_v - p_2)}}
        = (-1)^{p_2 - 2}.
    \end{align*}
    综上,
    \begin{equation*}
        s(p_1, p_2, p_3, \dots, p_n)
        = (-1)^{p_2 - 2}\, s(p_3, \dots, p_n).
        \qedhere
    \end{equation*}
\end{proof}

现在, 我们证明, ``qfaffian'' 就是 pfaffian.

\begin{theorem}
    设 \(A\) 是 \(2m\)~级反称阵.
    则
    % \(\operatorname{pf} {(A)} = \operatorname{qf} {(A)}\).
    % 所以,
    \begin{align*}
        \operatorname{pf} {(A)}
        = \sum_{\substack{
        1 \leq i_1, j_1, \dots, i_m, j_m \leq 2m; \\
        i_1, j_1, \dots, i_m, j_m\,\text{互不相同};   \\
        i_1 < j_1, \dots, i_m < j_m;              \\
                i_1 < \dots < i_m
            }}
        {s(i_1, j_1, \dots, i_m, j_m)\,
            [A]_{i_1,j_1} \dots [A]_{i_m,j_m}}.
    \end{align*}
\end{theorem}

\begin{proof}
    作命题 \(P(m)\):
    对任何 \(2m\)~级反称阵 \(A\),
    \(\operatorname{pf} {(A)} = \operatorname{qf} {(A)}\).
    我们用数学归纳法证明,
    对任何正整数 \(m\),
    \(P(m)\) 是对的.

    \(P(1)\) 是对的.

    我们设 \(P(m-1)\) 是对的.
    我们要由此证明, \(P(m)\) 是对的.

    任取 \(2m\)~级反称阵 \(A\).
    记 \(i_1 = 1\).
    则
    \begin{align*}
        \operatorname{pf} {(A)}
        = {} &
        \sum_{2 \leq j_1 \leq 2m}
        {(-1)^{j_1} [A]_{i_1,j_1}
        \operatorname{pf} {(A({i_1,j_1}|{i_1,j_1}))}}
        \\
        = {} &
        \sum_{\substack{
        i_1 = 1; \\
            i_1 < j_1 \leq 2m
        }}
        {(-1)^{j_1} [A]_{i_1,j_1}
        \operatorname{pf} {(A({i_1,j_1}|{i_1,j_1}))}}.
    \end{align*}
    由假定,
    \begin{align*}
        \operatorname{pf} {(A({i_1,j_1}|{i_1,j_1}))}
        = {} &
        \operatorname{qf} {(A({i_1,j_1}|{i_1,j_1}))}
        \\
        = {} &
        \sum_{\substack{
        1 \leq i_2, j_2, \dots, i_m, j_m \leq 2m; \\
        i_2, j_2, \dots, i_m, j_m \neq i_1;       \\
        i_2, j_2, \dots, i_m, j_m \neq j_1;       \\
        i_2, j_2, \dots, i_m, j_m\,\text{互不相同};   \\
        i_2 < j_2, \dots, i_m < j_m;              \\
                i_2 < \dots < i_m
            }}
        {s(i_2, j_2, \dots, i_m, j_m)\,
            [A]_{i_2,j_2} \dots [A]_{i_m,j_m}}.
    \end{align*}
    则
    \begin{align*}
             &
        \operatorname{pf} {(A)}
        \\
        = {} &
        \sum_{\substack{
        i_1 = 1;                                            \\
                i_1 < j_1 \leq 2m
            }}
        (-1)^{j_1} [A]_{i_1,j_1}
        \sum_{\substack{
        1 \leq i_2, j_2, \dots, i_m, j_m \leq 2m;           \\
        i_2, j_2, \dots, i_m, j_m \neq i_1;                 \\
        i_2, j_2, \dots, i_m, j_m \neq j_1;                 \\
        i_2, j_2, \dots, i_m, j_m\,\text{互不相同};             \\
        i_2 < j_2, \dots, i_m < j_m;                        \\
                i_2 < \dots < i_m
            }}
        {s(i_2, j_2, \dots, i_m, j_m)\,
            [A]_{i_2,j_2} \dots [A]_{i_m,j_m}}
        \\
        = {} &
        \sum_{\substack{
        i_1 = 1;                                            \\
                i_1 < j_1 \leq 2m
            }}
        \sum_{\substack{
        1 \leq i_2, j_2, \dots, i_m, j_m \leq 2m;           \\
        i_2, j_2, \dots, i_m, j_m \neq i_1;                 \\
        i_2, j_2, \dots, i_m, j_m \neq j_1;                 \\
        i_2, j_2, \dots, i_m, j_m\,\text{互不相同};             \\
        i_2 < j_2, \dots, i_m < j_m;                        \\
                i_2 < \dots < i_m
            }}
        { (-1)^{j_1} [A]_{i_1,j_1}
            s(i_2, j_2, \dots, i_m, j_m)\,
            [A]_{i_2,j_2} \dots [A]_{i_m,j_m}}
        \\
        = {} &
        \sum_{\substack{
        1 \leq i_1, j_1, i_2, j_2, \dots, i_m, j_m \leq 2m; \\
        i_1, j_1, i_2, j_2, \dots, i_m, j_m\,\text{互不相同};   \\
        i_1 < j_1, i_2 < j_2, \dots, i_m < j_m;             \\
                i_1 < i_2 < \dots < i_m
            }}
        { (-1)^{j_1-2}
            s(i_2, j_2, \dots, i_m, j_m)\,
            [A]_{i_1,j_1} [A]_{i_2,j_2} \dots [A]_{i_m,j_m}}
        \\
        = {} &
        \sum_{\substack{
        1 \leq i_1, j_1, \dots, i_m, j_m \leq 2m;           \\
        i_1, j_1, \dots, i_m, j_m\,\text{互不相同};             \\
        i_1 < j_1, \dots, i_m < j_m;                        \\
                i_1 < \dots < i_m
            }}
        {s(i_1, j_1, i_2, j_2, \dots, i_m, j_m)\,
            [A]_{i_1,j_1} [A]_{i_2,j_2} \dots [A]_{i_m,j_m}}
        \\
        = {} &
        \operatorname{qf} {(A)}.
    \end{align*}

    所以, \(P(m)\) 是对的.
    由数学归纳法, 待证命题成立.
\end{proof}

既然我们有了 pfaffian 的完全展开,
我们就不再说 ``qfaffian'',
且不再用记号 \(\operatorname{qf}\).

值得提, \(2m\)~级反称阵的 pfaffian 有
\((2m-1) \cdot (2m-3) \cdot \dots \cdot 1\) 项.
毕竟, 若互不相同的
\(i_1\), \(j_1\), \(\dots\), \(i_m\), \(j_m\)
适合
\(i_1 < j_1\), \(\dots\), \(i_m < j_m\),
且 \(i_1 < \dots < i_m\),
则 \(i_1\) 一定是 \(1\), \(2\), \(\dots\), \(2m\) 的最小的数, \(1\).
则 \(j_1\) 有 \(2m-1\) 种选择.
则 \(i_2\) 一定是剩下的 \(2m-2\) 个数的最小的数.
则 \(j_2\) 有 \(2m-3\) 种选择.
……
则 \(i_m\) 一定是剩下的 \(2\) 个数的较小的数.
则 \(j_m\) 有 \(1\) 种选择.
由分步乘法计数原理, 这样的
\(i_1\), \(j_1\), \(\dots\), \(i_m\), \(j_m\)
的个数是
\begin{align*}
    1 \cdot (2m-1) \cdot 1 \cdot (2m-3) \cdot \dots \cdot
    1 \cdot 1
    = (2m-1) \cdot (2m-3) \cdot \dots \cdot 1.
\end{align*}

Pfaffian 的完全展开有如下的变体:

\begin{theorem}
    设 \(A\) 是 \(2m\)~级反称阵.
    则
    \begin{align*}
             &
        \sum_{\substack{
        1 \leq u_1, v_1, \dots, u_m, v_m \leq 2m; \\
                u_1, v_1, \dots, u_m, v_m\,\text{互不相同}
            }}
        {s(u_1, v_1, \dots, u_m, v_m)\,
            [A]_{u_1,v_1} \dots [A]_{u_m,v_m}}
        \\
        = {} &
        2^m m!
        \sum_{\substack{
        1 \leq i_1, j_1, \dots, i_m, j_m \leq 2m; \\
        i_1, j_1, \dots, i_m, j_m\,\text{互不相同};   \\
        i_1 < j_1, \dots, i_m < j_m;              \\
                i_1 < \dots < i_m
            }}
        {s(i_1, j_1, \dots, i_m, j_m)\,
            [A]_{i_1,j_1} \dots [A]_{i_m,j_m}},
    \end{align*}
    其中,
    \(m! = m \cdot (m-1) \cdot \dots \cdot 1\).
    所以,
    \begin{align*}
        \operatorname{pf} {(A)} =
        \frac{1}{2^m m!}
        \sum_{\substack{
        1 \leq u_1, v_1, \dots, u_m, v_m \leq 2m; \\
                u_1, v_1, \dots, u_m, v_m\,\text{互不相同}
            }}
        {s(u_1, v_1, \dots, u_m, v_m)\,
            [A]_{u_1,v_1} \dots [A]_{u_m,v_m}}.
    \end{align*}
\end{theorem}

\begin{proof}
    数的乘法适合结合律与交换律.
    那么, 对任何项
    \begin{align*}
        s(u_1, v_1, \dots, u_m, v_m)\, [A]_{u_1,v_1} \dots [A]_{u_m,v_m},
    \end{align*}
    适当地交换
    \([A]_{u_1,v_1}\), \(\dots\), \([A]_{u_m,v_m}\)
    的次序,
    并注意, \([A]_{i,j} = -[A]_{j,i}\),
    我们可变它为 \(\sigma(u_1, v_1, \dots, u_m, v_m)\,
    [A]_{i_1,j_1} \dots [A]_{i_m,j_m}\),
    其中,
    \(\sigma(u_1, v_1, \dots, u_m, v_m)\)
    是由 \(u_1\), \(v_1\), \(\dots\), \(u_m\), \(v_m\) 确定的数
    (\(1\) 或 \(-1\)),
    \(i_1 < j_1\), \(\dots\), \(i_m < j_m\),
    且 \(i_1 < \dots < i_m\).
    我们说明,
    \begin{align*}
        \sigma(u_1, v_1, \dots, u_m, v_m) = s(i_1, j_1, \dots, i_m, j_m),
    \end{align*}
    从而
    \begin{align*}
        s(u_1, v_1, \dots, u_m, v_m)\, [A]_{u_1,v_1} \dots [A]_{u_m,v_m}
        =
        s(i_1, j_1, \dots, i_m, j_m)\, [A]_{i_1,j_1} \dots [A]_{i_m,j_m}.
    \end{align*}

    回想, 交换排列的二个数, 则它的符号变为原符号的相反数
    (见 ``排列'').
    那么,
    \begin{align*}
        s(\dots, u_p, v_p, \dots, u_q, v_q, \dots)
        = {} &
        (-1)\, s(\dots, u_q, v_p, \dots, u_p, v_q, \dots)
        \\
        = {} &
        (-1)\, (-1)\, s(\dots, u_q, v_q, \dots, u_p, v_p, \dots)
        \\
        = {} &
        s(\dots, u_q, v_q, \dots, u_p, v_p, \dots).
    \end{align*}
    这里, 未写的元是不变的, 下同.
    则
    \begin{align*}
             &
        s(\dots, u_p, v_p, \dots, u_q, v_q, \dots)\,
        (\dots [A]_{u_p,v_p} \dots [A]_{u_q,v_q} \dots)
        \\
        = {} &
        s(\dots, u_q, v_q, \dots, u_p, v_p, \dots)\,
        (\dots [A]_{u_q,v_q} \dots [A]_{u_p,v_p} \dots).
    \end{align*}
    另外,
    \begin{align*}
        s(\dots, u_p, v_p, \dots)
        = (-1)\, s(\dots, u_q, v_p, \dots),
    \end{align*}
    且
    \begin{align*}
        (\dots [A]_{u_q,v_q} \dots)
        = (-1)\, (\dots [A]_{v_q,u_q} \dots),
    \end{align*}
    故
    \begin{align*}
        s(\dots, u_p, v_p, \dots)\, (\dots [A]_{u_q,v_q} \dots)
        = s(\dots, v_p, u_p, \dots)\, (\dots [A]_{v_q,u_q} \dots).
    \end{align*}

    对 \(1\), \(2\), \(\dots\), \(2m\) 的任何排列
    \(u_1\), \(v_1\), \(\dots\), \(u_m\), \(v_m\),
    我们分它的 \(2m\) 个数为 \(m\) 个偶
    \((u_1, v_1)\), \(\dots\), \((u_m, v_m)\).
    考虑由此排列确定的项
    \begin{align*}
        s(u_1, v_1, \dots, u_m, v_m)\, [A]_{u_1,v_1} \dots [A]_{u_m,v_m}.
    \end{align*}
    前面的计算说明,
    交换这 \(m\) 个偶的次序不改变项的结果,
    且交换一个偶的二个数的次序也不改变项的结果.
    故
    \begin{align*}
        s(u_1, v_1, \dots, u_m, v_m)\, [A]_{u_1,v_1} \dots [A]_{u_m,v_m}
        =
        s(i_1, j_1, \dots, i_m, j_m)\, [A]_{i_1,j_1} \dots [A]_{i_m,j_m},
    \end{align*}
    其中, \(i_1 < j_1\), \(\dots\), \(i_m < j_m\),
    且 \(i_1 < \dots < i_m\).

    注意,
    \(1\), \(2\), \(\dots\), \(2m\) 的任何排列
    \(u_1\), \(v_1\), \(\dots\), \(u_m\), \(v_m\)
    可被变为某个
    \(i_1\), \(j_1\), \(\dots\), \(i_m\), \(j_m\),
    其中, \(i_1 < j_1\), \(\dots\), \(i_m < j_m\),
    且 \(i_1 < \dots < i_m\).
    反过来,
    对于任何这样的
    \(i_1\), \(j_1\), \(\dots\), \(i_m\), \(j_m\),
    我们分它的 \(2m\) 个数为 \(m\) 个偶
    \((i_1, j_1)\), \(\dots\), \((i_m, j_m)\),
    并自由地交换偶的次序,
    与一个偶的二个数的次序,
    可得
    \(1\), \(2\), \(\dots\), \(2m\) 的一个排列.
    我们有
    \(m \cdot (m-1) \cdot \dots \cdot 1 = m!\)
    种交换偶的次序的方式.
    偶的次序被给定后,
    我们有
    \(2 \cdot 2 \cdot \dots \cdot 2 \text{ (\(m\) 个 \(2\))}
    = 2^m\) 种交换偶的二个数的次序的方式.
    所以, 一个
    \(i_1\), \(j_1\), \(\dots\), \(i_m\), \(j_m\)
    (\(i_1 < j_1\), \(\dots\), \(i_m < j_m\),
    且 \(i_1 < \dots < i_m\))
    给出 \(2^m m!\) 个
    \(1\), \(2\), \(\dots\), \(2m\) 的排列,
    且由这些排列确定的项是相等的.
    故
    \begin{align*}
             &
        \sum_{\substack{
        1 \leq u_1, v_1, \dots, u_m, v_m \leq 2m; \\
                u_1, v_1, \dots, u_m, v_m\,\text{互不相同}
            }}
        {s(u_1, v_1, \dots, u_m, v_m)\,
            [A]_{u_1,v_1} \dots [A]_{u_m,v_m}}
        \\
        = {} &
        2^m m!
        \sum_{\substack{
        1 \leq i_1, j_1, \dots, i_m, j_m \leq 2m; \\
        i_1, j_1, \dots, i_m, j_m\,\text{互不相同};   \\
        i_1 < j_1, \dots, i_m < j_m;              \\
                i_1 < \dots < i_m
            }}
        {s(i_1, j_1, \dots, i_m, j_m)\,
            [A]_{i_1,j_1} \dots [A]_{i_m,j_m}}.
    \end{align*}
\end{proof}

最后, 我们用 pfaffian 的完全展开证明 pfaffian 的性质.

% a manual pagebreak
\clearpage

\begin{theorem}
    设 \(A\) 是 \(\ell\)~级反称阵.
    设 \(P\) 是 \(n \times \ell\)~阵.

    (1)
    若 \(n > \ell\),
    则 \(\operatorname{pf} {(P A P^{\mathrm{T}})} = 0\).

    (2)
    若 \(n \leq \ell\),
    则
    \begin{align*}
             &
        \operatorname{pf} {(P A P^{\mathrm{T}})}
        \\
        = {} &
        \sum_{1 \leq j_1 < j_2 < \dots < j_{n} \leq \ell}
        {\det {\left(
                P\binom{1,2,\dots,n}{j_1,j_2,\dots,j_{n}}
                \right)}
            \operatorname{pf} {\left(
                A\binom{j_1,j_2,\dots,j_{n}}{j_1,j_2,\dots,j_{n}}
                \right)}}.
    \end{align*}
    特别地, 若 \(n = \ell\),
    则因适合条件
    \(1 \leq j_1 < j_2 < \dots < j_{n} \leq \ell\)
    的
    \(j_1\), \(j_2\), \(\dots\), \(j_{n}\),
    是, 且只能是,
    \(1\), \(2\), \(\dots\), \(n\),
    \begin{align*}
        \operatorname{pf} {(P A P^{\mathrm{T}})}
        = {} &
        \det {\left(
            P\binom{1,2,\dots,n}{1,2,\dots,n}
            \right)}
        \operatorname{pf} {\left(
            A\binom{1,2,\dots,n}{1,2,\dots,n}
            \right)}
        \\
        = {} &
        \det {(P)} \operatorname{pf} {(A)}.
    \end{align*}
\end{theorem}

\begin{proof}
    若 \(n\) 是奇数, 则
    \(P A P^{\mathrm{T}}\)
    与
    \(\displaystyle
    A\binom{j_1,j_2,\dots,j_{n}}{j_1,j_2,\dots,j_{n}}\)
    都是奇数级反称阵,
    故它们的 pfaffian 都是 \(0\).
    则这是显然的.

    以下, 设 \(n = 2m\) 是偶数.

    注意,
    \begin{align*}
        [P A P^{\mathrm{T}}]_{i,j}
        = {} &
        \sum_{1 \leq t \leq \ell}
        {[P A]_{i,t} [P^{\mathrm{T}}]_{t,j}}
        \\
        = {} &
        \sum_{1 \leq t \leq \ell}
        {
            \left(
            \sum_{1 \leq r \leq \ell}
            {[P]_{i,r} [A]_{r,t}}
            \right)
            [P]_{j,t}
        }
        \\
        = {} &
        \sum_{1 \leq r, t \leq \ell}
        {[P]_{i,r} [P]_{j,t} [A]_{r,t}}.
    \end{align*}
    为方便, 我们简单地写
    ``\(1 \leq u_1, v_1, \dots, u_m, v_m \leq x\)''
    为 ``\(u, v \text{:}\, m \text{:}\, x\)'',
    且简单地写 ``\(u_1, v_1, \dots, u_m, v_m\) 互不相同''
    为 ``\(u, v \,\text{D}\)''.
    则
    \begin{align*}
             &
        \operatorname{pf} {(P A P^{\mathrm{T}})}
        \\
        = {} &
        \frac{1}{2^m m!}
        \sum_{\substack{
        u, v \text{:}\, m \text{:}\, 2m; \\
                u, v \,\text{D}
            }}
        {s(u_1, v_1, \dots, u_m, v_m)\,
            [P A P^{\mathrm{T}}]_{u_1,v_1}
            \dots
            [P A P^{\mathrm{T}}]_{u_m,v_m}}
        \\
        = {} &
        \frac{1}{2^m m!}
        \sum_{\substack{
        u, v \text{:}\, m \text{:}\, 2m; \\
                u, v \,\text{D}
            }}
        s(u_1, v_1, \dots, u_m, v_m)
        \left(
        \sum_{1 \leq r_1, t_1 \leq \ell}
        {[P]_{u_1,r_1} [P]_{v_1,t_1} [A]_{r_1,t_1}}
        \right)
        \\
             &
        \quad \quad \quad \quad \quad \,
        {} \cdot
        \dots
        \cdot
        \left(
        \sum_{1 \leq r_m, t_m \leq \ell}
        {[P]_{u_m,r_m} [P]_{v_m,t_m} [A]_{r_m,t_m}}
        \right)
        \\
        = {} &
        \frac{1}{2^m m!}
        \sum_{\substack{
        u, v \text{:}\, m \text{:}\, 2m; \\
                u, v \,\text{D}
            }}
        s(u_1, v_1, \dots, u_m, v_m)     \\
             &
        \quad \quad \quad \quad \quad
        {} \cdot
        \sum_{r, t \text{:}\, m \text{:}\, \ell}
        [P]_{u_1,r_1} [P]_{v_1,t_1} \dots [P]_{u_m,r_m} [P]_{v_m,t_m}
        \cdot [A]_{r_1,t_1} \dots [A]_{r_m,t_m}
        \\
        = {} &
        \frac{1}{2^m m!}
        \sum_{\substack{
        u, v \text{:}\, m \text{:}\, 2m; \\
                u, v \,\text{D}
            }}
        \sum_{r, t \text{:}\, m \text{:}\, \ell}
        s(u_1, v_1, \dots, u_m, v_m)
        \,
        [P]_{u_1,r_1} [P]_{v_1,t_1} \dots [P]_{u_m,r_m} [P]_{v_m,t_m}
        \\
             &
        \quad \quad \quad \quad
        \quad \quad \quad \quad
        {} \cdot [A]_{r_1,t_1} \dots [A]_{r_m,t_m}
        \\
        = {} &
        \frac{1}{2^m m!}
        \sum_{r, t \text{:}\, m \text{:}\, \ell}
        \sum_{\substack{
        u, v \text{:}\, m \text{:}\, 2m; \\
                u, v \,\text{D}
            }}
        s(u_1, v_1, \dots, u_m, v_m)
        \,
        [P]_{u_1,r_1} [P]_{v_1,t_1} \dots [P]_{u_m,r_m} [P]_{v_m,t_m}
        \\
             &
        \quad \quad \quad \quad
        \quad \quad \quad \quad
        {} \cdot [A]_{r_1,t_1} \dots [A]_{r_m,t_m}
        \\
        = {} &
        \frac{1}{2^m m!}
        \sum_{r, t \text{:}\, m \text{:}\, \ell}
        \left(
        \sum_{\substack{
        u, v \text{:}\, m \text{:}\, 2m; \\
                u, v \,\text{D}
            }}
        s(u_1, v_1, \dots, u_m, v_m)
        \,
        [P]_{u_1,r_1} [P]_{v_1,t_1} \dots [P]_{u_m,r_m} [P]_{v_m,t_m}
        \right)
        \\
             &
        \quad \quad \quad \quad \quad
        {} \cdot [A]_{r_1,t_1} \dots [A]_{r_m,t_m}.
    \end{align*}
    作 \(2m\)~级阵 \(Q(r_1, t_1, \dots, r_m, t_m)\) 如下:
    \begin{align*}
        [Q(r_1, t_1, \dots, r_m, t_m)]_{i,2j-1} & = [P]_{i,r_j}, \\
        [Q(r_1, t_1, \dots, r_m, t_m)]_{i,2j}   & = [P]_{i,t_j}.
    \end{align*}
    则
    \begin{align*}
             &
        \sum_{\substack{
        u, v \text{:}\, m \text{:}\, 2m; \\
                u, v \,\text{D}
            }}
        s(u_1, v_1, \dots, u_m, v_m)
        \,
        [P]_{u_1,r_1} [P]_{v_1,t_1} \dots [P]_{u_m,r_m} [P]_{v_m,t_m}
        \\
        = {} &
        \det {(Q(r_1, t_1, \dots, r_m, t_m))}.
    \end{align*}
    则
    \begin{align*}
        \operatorname{pf} {(P A P^{\mathrm{T}})}
        = {} &
        \frac{1}{2^m m!}
        \sum_{r, t \text{:}\, m \text{:}\, \ell}
        \det {(Q(r_1, t_1, \dots, r_m, t_m))}
        \,
        [A]_{r_1,t_1} \dots [A]_{r_m,t_m}.
    \end{align*}
    若在 \(r_1\), \(t_1\), \(\dots\), \(r_m\), \(t_m\),
    有二个数相同,
    则 \(Q(r_1, t_1, \dots, r_m, t_m)\) 有二列相同.
    则 \(\det {(Q(r_1, t_1, \dots, r_m, t_m))} = 0\).

    (1)
    若 \(2m > \ell\),
    则在 \(r_1\), \(t_1\), \(\dots\), \(r_m\), \(t_m\),
    总有二个数相同.
    则和的每一项都是 \(0\).
    故 \(\operatorname{pf} {(P A P^{\mathrm{T}})} = 0\).

    (2)
    若 \(2m \leq \ell\),
    则 (见第一章, 节~\malneprasekcio{30})
    \begin{align*}
             &
        \operatorname{pf} {(P A P^{\mathrm{T}})}
        \\
        = {} &
        \frac{1}{2^m m!}
        \sum_{1 \leq j_1 < \dots < j_{2m} \leq \ell}
        \sum_{\substack{
        r_1, t_1, \dots, r_m, t_m\,\text{是} \\
                j_1, \dots, j_{2m}\,\text{的排列}
            }}
        \det {(Q(r_1, t_1, \dots, r_m, t_m))}
        \,
        [A]_{r_1,t_1} \dots [A]_{r_m,t_m}.
    \end{align*}
    因为 \(j_1 < j_2 < \dots < j_{2m}\),
    且 \(r_1\), \(t_1\), \(\dots\), \(r_m\), \(t_m\)
    是 \(j_1\), \(j_2\), \(\dots\), \(j_{2m}\) 的排列
    (见第一章, 节~\malneprasekcio{30})
    \begin{align*}
             &
        \det {(Q(r_1, t_1, \dots, r_m, t_m))}
        \\
        = {} &
        \det {(Q(j_1, j_2, \dots, j_{2m}))}\,
        s(r_1, t_1, \dots, r_m, t_m)
        \\
        = {} &
        \det {\left(
            P\binom{1,2,\dots,2m}{j_1,j_2,\dots,j_{2m}}
            \right)}\,
        s(r_1, t_1, \dots, r_m, t_m).
    \end{align*}
    则
    \begin{align*}
             &
        \operatorname{pf} {(P A P^{\mathrm{T}})}
        \\
        = {} &
        \frac{1}{2^m m!}
        \sum_{1 \leq j_1 < \dots < j_{2m} \leq \ell}
        \sum_{\substack{
        r_1, t_1, \dots, r_m, t_m\,\text{是} \\
                j_1, \dots, j_{2m}\,\text{的排列}
            }}
        \det {\left(
            P\binom{1,\dots,2m}{j_1,\dots,j_{2m}}
            \right)}\,
        s(r_1, t_1, \dots, r_m, t_m)
        \\
             &
        \quad \quad \quad \quad \quad \quad \quad \quad \quad \quad \quad \quad \quad
        \cdot
        [A]_{r_1,t_1} \dots [A]_{r_m,t_m}
        \\
        = {} &
        \frac{1}{2^m m!}
        \sum_{1 \leq j_1 < \dots < j_{2m} \leq \ell}
        \det {\left(
            P\binom{1,\dots,2m}{j_1,\dots,j_{2m}}
            \right)}
        \\
             &
        \quad \quad \quad \quad \quad \quad \quad \quad
        \cdot
        \sum_{\substack{
        r_1, t_1, \dots, r_m, t_m\,\text{是} \\
                j_1, \dots, j_{2m}\,\text{的排列}
            }}
        s(r_1, t_1, \dots, r_m, t_m)\,
        [A]_{r_1,t_1} \dots [A]_{r_m,t_m}
        \\
        = {} &
        \sum_{1 \leq j_1 < \dots < j_{2m} \leq \ell}
        \det {\left(
            P\binom{1,\dots,2m}{j_1,\dots,j_{2m}}
            \right)}
        \\
             &
        \quad \quad \quad \quad \quad
        \cdot
        \frac{1}{2^m m!}
        \sum_{\substack{
        r_1, t_1, \dots, r_m, t_m\,\text{是} \\
                j_1, \dots, j_{2m}\,\text{的排列}
            }}
        s(r_1, t_1, \dots, r_m, t_m)\,
        [A]_{r_1,t_1} \dots [A]_{r_m,t_m}.
    \end{align*}

    设
    \begin{align*}
        H(j_1, j_2, \dots, j_{2m}) =
        A\binom{j_1,j_2,\dots,j_{2m}}{j_1,j_2,\dots,j_{2m}}.
    \end{align*}
    不难看出,
    \([H(j_1, j_2, \dots, j_{2m})]_{u,v} = [A]_{j_u,j_v}\).
    记 \(f(j_a) = a\)
    (\(a = 1\), \(2\), \(\dots\), \(2m\)).
    设 \(g_1\), \(g_2\), \(\dots\), \(g_{2m}\)
    是 \(j_1\), \(j_2\), \(\dots\), \(j_{2m}\)
    的一个排列.
    那么, \(f(g_1)\), \(f(g_2)\), \(\dots\), \(f(g_{2m})\)
    是 \(1\), \(2\), \(\dots\), \(2m\) 的一个排列.
    反过来, 若 \(w_1\), \(w_2\), \(\dots\), \(w_{2m}\)
    是 \(1\), \(2\), \(\dots\), \(2m\) 的一个排列,
    则 \(j_{w_1}\), \(j_{w_2}\), \(\dots\), \(j_{w_{2m}}\)
    是 \(j_1\), \(j_2\), \(\dots\), \(j_{2m}\) 的一个排列.
    并且, 若 \(g_a < g_b\), 必有 \(f(g_a) < f(g_b)\).
    反过来, 若 \(f(g_a) < f(g_b)\), 必有 \(g_a < g_b\).
    从而
    \(\operatorname{sgn} {(g_b - g_a)}
    = \operatorname{sgn} {(f(g_b) - f(g_a))}\).
    则
    \(s(g_1, g_2, \dots, g_{2m})
    = s(f(g_1), f(g_2), \dots, f(g_{2m}))\).
    故
    \begin{align*}
             &
        \frac{1}{2^m m!}
        \sum_{\substack{
        r_1, t_1, \dots, r_m, t_m\,\text{是} \\
                j_1, \dots, j_{2m}\,\text{的排列}
            }}
        s(r_1, t_1, \dots, r_m, t_m)\,
        [A]_{r_1,t_1} \dots [A]_{r_m,t_m}
        \\
        = {} &
        \frac{1}{2^m m!}
        \sum_{\substack{
        r_1, t_1, \dots, r_m, t_m\,\text{是} \\
                j_1, \dots, j_{2m}\,\text{的排列}
            }}
        s(f(r_1), f(t_1), \dots, f(r_m), f(t_m))
        \\
             &
        \quad \quad \quad \quad \quad \quad \quad \quad
        \cdot
        [H(j_1, \dots, j_{2m})]_{f(r_1),f(t_1)}
        \dots
        [H(j_1, \dots, j_{2m})]_{f(r_m),f(t_m)}
        \\
        = {} &
        \operatorname{pf} {(H(j_1, \dots, j_{2m}))}
        \\
        = {} &
        \operatorname{pf} {\left(
            A\binom{j_1,\dots,j_{2m}}{j_1,\dots,j_{2m}}
            \right)}.
    \end{align*}
    从而
    \begin{align*}
        \operatorname{pf} {(P A P^{\mathrm{T}})}
        = {} &
        \sum_{1 \leq j_1 < \dots < j_{2m} \leq \ell}
        {\det {\left(
                P\binom{1,\dots,2m}{j_1,\dots,j_{2m}}
                \right)}
            \operatorname{pf} {\left(
                A\binom{j_1,\dots,j_{2m}}{j_1,\dots,j_{2m}}
                \right)}}.
    \end{align*}
\end{proof}

% a manual pagebreak
\clearpage

\begin{theorem}
    设 \(A\) 是 \(n\)~级反称阵.
    则 \(\det {(A)}
    = (\operatorname{pf} {(A)})^2\).
\end{theorem}

\begin{proof}
    设 \(n\) 是奇数.
    则 \(\det {(A)} = 0\),
    且 \(\operatorname{pf} {(A)} = 0\).

    再设 \(n = 2m\) 是偶数.
    我们知道,
    \begin{align*}
        \operatorname{pf} {(P A P^{\mathrm{T}})}
        = \det {(P)} \operatorname{pf} {(A)}.
    \end{align*}
    我们取 \(P\) 为
    \(2m\)~级倍加阵 \(E(2m; p, q; s)\),
    其中, \(p \neq q\), 且
    \begin{align*}
        [E(2m; p, q; s)]_{i,j}
        = \begin{cases}
              s,              & \text{\(i = p\), 且 \(j = q\)}; \\
              [I_{2m}]_{i,j}, & \text{其他},
          \end{cases}
    \end{align*}
    则我们知道,
    作一次列的倍加与一次对应的行的倍加不改变 pfaffian,
    因为倍加阵的行列式是 \(1\).

    利用若干次列的倍加, 与对应的行的倍加
    (先列后行, 交替地作),
    我们可变 \(A\) 为
    \begin{align*}
        B =
        \begin{bmatrix}
            0      & b_1    & 0      & 0      & \cdots & 0      & 0      \\
            -b_1   & 0      & 0      & 0      & \cdots & 0      & 0      \\
            0      & 0      & 0      & b_2    & \cdots & 0      & 0      \\
            0      & 0      & -b_2   & 0      & \cdots & 0      & 0      \\
            \vdots & \vdots & \vdots & \vdots & \ddots & \vdots & \vdots \\
            0      & 0      & 0      & 0      & \cdots & 0      & b_m    \\
            0      & 0      & 0      & 0      & \cdots & -b_m   & 0      \\
        \end{bmatrix}.
    \end{align*}
    则 \(\det {(B)} = \det {(A)}\),
    且 \(\operatorname{pf} {(B)} = \operatorname{pf} {(A)}\).
    最后, 注意, \(\det {(B)}
    = (\operatorname{pf} {(B)})^2\),
    故 \(\det {(A)}
    = (\operatorname{pf} {(A)})^2\).
\end{proof}

\end{document}
